(1) 小角摆的回复力为 $-mg(x,y)/\ell$，固有频率 $\omega_0=\sqrt{g/\ell}$；地球自转矢量在当地竖直方向的分量为 $f=\Omega_\oplus\sin\varphi$。在 $\Omega_\oplus\ll\omega_0$ 且水平尺寸远小于地球半径的假设下，水平自转分量主要带来更高阶竖直耦合，可在本题一阶模型略去。
(2) 旋转坐标系科里奥利加速度为 $-2\boldsymbol\Omega_\oplus\times\dot{\boldsymbol r}$；采用东、北、上的右手坐标，得到 $\ddot x-2f\dot y+\omega_0^2x=0$ 与 $\ddot y+2f\dot x+\omega_0^2y=0$。两方程的符号成对相反，是因为科里奥利加速度始终垂直于瞬时水平速度。
(3) 令 $z=x+iy$，方程合并为 $\ddot z+2if\dot z+\omega_0^2z=0$。用 $z\propto e^{-i\omega t}$ 得 $\omega=f\pm\sqrt{\omega_0^2+f^2}$；两圆偏振分量的正频率大小为 $\sqrt{\omega_0^2+f^2}\pm f$。同线偏振初态可表示为这两个反向转动模叠加，地面轨道主轴以 $-f=-\Omega_\oplus\sin\varphi$ 的角速度旋转，在北半球从上往下看为顺时针；一阶进动周期为 $2\pi/|f|$。
(4) 将方程分别乘以 $m\dot x,m\dot y$ 并相加，交叉项 $-2mf\dot x\dot y+2mf\dot y\dot x$ 抵消，所以 $d[m(\dot x^2+\dot y^2+\omega_0^2x^2+\omega_0^2y^2)/2]/dt=0$。赤道 $\sin\varphi=0$，摆面不发生一阶傅科进动；两极 $|f|=\Omega_\oplus$，一恒星日转一周，方位方向依半球而异。
